% =====================================================================
%  The Irreducible Boundary of Bounded Phase Space
%  A self-contained topological theory of division, boundary thickness,
%  and the resolution floor, with physical realisation.
%
%  Self-contained. No reliance on any unpublished manuscript.
% =====================================================================
\documentclass[11pt,a4paper]{article}

\usepackage[utf8]{inputenc}
\usepackage[T1]{fontenc}
\usepackage{lmodern}
\usepackage[a4paper,margin=2.5cm]{geometry}
\usepackage{amsmath,amssymb,amsthm,mathtools}
\usepackage{bm}
\usepackage{mathrsfs}
\usepackage{booktabs}
\usepackage{array}
\usepackage{enumitem}
\usepackage{microtype}
\usepackage[round,sort&compress]{natbib}
\usepackage{tikz}
\usetikzlibrary{arrows.meta,positioning,calc,decorations.pathmorphing,fit}
\usepackage[colorlinks=true,linkcolor=blue!45!black,
            citecolor=blue!45!black,urlcolor=blue!45!black]{hyperref}
\usepackage{cleveref}

% ---- Theorem environments ----
\newtheoremstyle{thmstyle}{\topsep}{\topsep}{\itshape}{}{\bfseries}{.}{.5em}{}
\theoremstyle{thmstyle}
\newtheorem{theorem}{Theorem}[section]
\newtheorem{lemma}[theorem]{Lemma}
\newtheorem{proposition}[theorem]{Proposition}
\newtheorem{corollary}[theorem]{Corollary}
\theoremstyle{definition}
\newtheorem{definition}[theorem]{Definition}
\newtheorem{axiom}{Axiom}
\newtheorem{principle}[theorem]{Principle}
\newtheorem{example}[theorem]{Example}
\newtheorem{construction}[theorem]{Construction}
\theoremstyle{remark}
\newtheorem{remark}[theorem]{Remark}
\newtheorem{notation}[theorem]{Notation}

\crefname{axiom}{Axiom}{Axioms}
\Crefname{axiom}{Axiom}{Axioms}
\crefname{principle}{Principle}{Principles}
\Crefname{principle}{Principle}{Principles}

% ---- Shorthand ----
\newcommand{\R}{\mathbb{R}}
\newcommand{\N}{\mathbb{N}}
\newcommand{\Z}{\mathbb{Z}}
\newcommand{\Phase}{\Omega}
\newcommand{\Bdy}{\partial}
\newcommand{\interior}{\mathrm{int}}
\newcommand{\cl}{\mathrm{cl}}
\newcommand{\diam}{\mathrm{diam}}
\newcommand{\dist}{\mathrm{dist}}
\newcommand{\Leb}{\mathcal{L}}
\newcommand{\Haus}{\mathcal{H}}
\newcommand{\thick}{\beta}          % boundary thickness / floor
\newcommand{\reso}{\delta}          % resolution
\newcommand{\floor}{\thick}
\newcommand{\Part}{\mathscr{P}}
\newcommand{\cell}{\mathsf{c}}
\newcommand{\eps}{\varepsilon}
\newcommand{\kB}{k_{\mathrm{B}}}
\newcommand{\Sval}{S}
\newcommand{\receiver}{\mathcal{R}}
\newcommand{\KK}{\mathcal{K}}
\newcommand{\XX}{\mathcal{X}}
\newcommand{\norm}[1]{\lVert#1\rVert}
\newcommand{\abs}[1]{\lvert#1\rvert}
\newcommand{\dd}{\,\mathrm{d}}
\newcommand{\restr}{\!\restriction\!}

\title{\textbf{The Irreducible Boundary of Bounded Phase Space:\\[2pt]
Division, Boundary Thickness, and the Resolution Floor}}

\author{%
Kundai Farai Sachikonye\\
\small Technical University of Munich\\
\small Chair of Proteomics and Bioanalytics\\
\small \texttt{kundai.sachikonye@wzw.tum.de}}

\date{\today}

\begin{document}
\maketitle

% =====================================================================
\begin{abstract}
\noindent
We prove that to know a thing is to divide it from everything else, that
every such division requires a boundary of strictly positive thickness, and
that the contents of that thickness form an irreducible residue which no
bounded knower can eliminate. The development is topological and
measure-theoretic and rests on a single hypothesis---that the relevant state
space is bounded and finitely resolvable (the \emph{bounded phase space}
hypothesis). From it we derive a \emph{Boundary--Thickness Theorem}: in a
bounded space carrying a finite resolution, no two complementary regions can
be separated by a set of zero content; every separator carries positive
content $\thick>0$, the \emph{resolution floor}. We give three independent
derivations of the floor---a representational one (faithful images of a
consistent invariant cannot be the invariant), a thermodynamic one (the cost
of compressing a separator to zero content diverges, and a single binary
distinction costs at least $\kB T\ln 2$), and a cardinality one---and prove
they bound the \emph{same} constant. We prove a \emph{Detectability Theorem}:
a bounded object registers (``is distinguishable'') if and only if its scale
exceeds its own boundary thickness; sub-floor objects are engulfed by their
separator and cannot be individuated. The classical sorites paradox is then a
corollary, not a puzzle: there is no sharp grain at which a heap appears
because the heap--non-heap separator has positive thickness, and detection is
emergence from that thickness. We prove an \emph{Inquiry--Cessation Theorem}:
thinning a separator toward zero content has diverging cost and bounded
return, so a finite knower rationally halts at the floor, whence knowledge
attaches to bounded wholes and never to sub-floor constituents---a structure
we show is visible in ordinary attribution. Finally we exhibit the physical
realisation: the same hypothesis forces oscillatory dynamics, a discrete
partition spectrum, the coordinate quadruple $(n,\ell,m,s)$ with shell
capacity $2n^2$, an exclusion principle, the compression law, and a Fisher
embedding---so that the floor of the epistemic theory and the floor of the
physical theory are one constant. Every claim is proved from stated
definitions; uncertain material has been omitted.
\end{abstract}

\noindent\textbf{Keywords:} bounded phase space; topological boundary;
boundary thickness; resolution floor; sorites paradox; vagueness; partition;
compression cost; Landauer bound; information geometry; irreducible residue.

\tableofcontents

% =====================================================================
\section{Introduction}
\label{sec:intro}
% =====================================================================

\subsection{Thesis}

This paper defends, and proves under a single explicit hypothesis, the
following chain.

\begin{enumerate}[label=(\roman*),leftmargin=2.2em]
\item \emph{To know a thing is to divide it from everything else.} Knowledge
of an object $A$ is the capacity to tell $A$ from not-$A$; this is a partition
of the state space into $A$ and its complement.
\item \emph{Every division requires a boundary, and the boundary has positive
thickness.} In a bounded, finitely resolvable space there is no zero-content
separator between a region and its complement. The separator carries a strictly
positive content $\thick>0$.
\item \emph{The boundary's interior is an irreducible residue.} The contents of
the thickness belong determinately to neither $A$ nor not-$A$. They are the
operational form of the limitative ``residue'' familiar from logic and physics.
\item \emph{Detection is emergence from the boundary.} A bounded object is
distinguishable exactly when its own scale exceeds its boundary thickness.
Below that scale it is engulfed by its separator and cannot be individuated.
\item \emph{Inquiry rationally halts at the floor.} Thinning a separator toward
zero content costs without bound for bounded return; a finite knower stops, and
knowledge thereafter attaches to bounded wholes, never to their sub-floor parts.
\end{enumerate}

The single hypothesis is that the space in question is bounded and admits a
finite-resolution partition---the \emph{bounded phase space} hypothesis
(\Cref{ax:bps}). Everything else is theorem.

\subsection{Why ``thickness''}

The everyday and the technical meet at one point. A single grain of wheat
makes no sound when dropped; a heap makes a sound. Classically this is a
paradox: adding one grain never converts silence to sound, yet a heap sounds.
We resolve it by inverting the usual reading. The grain is silent not because
it is small in the absolute, but because its \emph{boundary is too thick
relative to it}: the layered separator one would have to resolve to isolate
``this grain'' from ``everything else'' is larger than the grain, so the grain
is buried inside its own boundary and never stands out as a figure against a
ground. The heap sounds because it is large compared with the same separator:
it clears its boundary and registers. There is no critical grain because
``registering'' is emergence across a band of positive width, not a crossing of
a point. The width of that band is the floor $\thick$. This is
\Cref{thm:detectability} below, and the sorites paradox is its
\Cref{cor:sorites}.

\subsection{Why ``bounded''}

Boundedness is not a convenience; it is the source of the floor. On an
unbounded space one could, in principle, separate a region from its complement
by an arbitrarily thin frontier and resolve structure to arbitrary depth at
fixed cost. It is precisely the finiteness of the accessible space---and the
finiteness of the resolution any embedded knower can purchase---that forbids the
zero-thickness cut. We make the hypothesis exactly strong enough to force the
floor and no stronger, and we mark, at each step, where boundedness is used.

\subsection{Two faces of one constant}

The floor $\thick$ has an epistemic face and a physical face, and a central aim
of the paper is to prove they coincide. Epistemically, $\thick$ is the smallest
distinction a bounded knower can draw---the thickness of its sharpest boundary.
Physically, the same hypothesis (\Cref{ax:bps}) forces a partition spectrum,
a compression law whose cost diverges as a cell is squeezed to zero measure
(\Cref{thm:compression}), and a minimum energy $\kB T\ln 2$ per binary
distinction (\Cref{cor:landauer}); these pin a physical floor. We prove
(\Cref{thm:floor-agreement}) that the epistemic and physical floors are not two
constants but bounds on one.

\subsection{Relation to existing treatments}

The sorites paradox and vagueness have large literatures: degree theories and
fuzzy logic~\citep{zadeh1965}, supervaluationism~\citep{fine1975,keefe2000},
epistemicism~\citep{williamson1994}, and the classical analyses
of~\citet{black1937}; see~\citet{huttegger2023} for a survey. Our account is
none of these: it is not a logic of partial truth, nor a claim that a sharp
boundary exists but is unknowable, but a \emph{geometric} statement that the
separator is a set of positive content whose interior is determinately
category-less. The residue connects structurally to limitative results in
logic~\citep{godel1931,turing1936} and to the thermodynamics of
computation~\citep{landauer1961,bennett1982}; we use these as analogues and, in
the physical sections, as tools, never as premises. The topological and
measure-theoretic machinery is standard~\citep{kelley1955,munkres2000,%
halmos1950,federer1969,mattila1995}, as is the ergodic and spectral
background~\citep{poincare1890,birkhoff1931,koopman1931,walters1982,%
reedsimon1980} and the information geometry~\citep{fisher1925,cencov1982,%
amari2016}.

\subsection{Method and discipline}

The paper is written to a single standard: every substantive statement is a
definition, a theorem proved here, or a citation of a standard result compared
against such a theorem. Where a tempting claim could not be proved at this
standard it has been removed rather than hedged. Physical interpretation is
confined to remarks and to \Cref{sec:physics}; the mathematical spine
(\Cref{sec:setting,sec:thickness,sec:floor,sec:detect,sec:residue,sec:inquiry})
stands without it.

\subsection{Organisation}

\Cref{sec:setting} fixes the topological and measure-theoretic setting and
states the bounded phase space hypothesis. \Cref{sec:thickness} proves the
Boundary--Thickness Theorem. \Cref{sec:floor} gives the three derivations of
the floor and proves their agreement. \Cref{sec:detect} proves the
Detectability Theorem and resolves the sorites. \Cref{sec:residue} develops
the topology of the residue (the boundary interior) and its non-attribution.
\Cref{sec:inquiry} proves the Inquiry--Cessation Theorem and the
bounded-attribution capstone. \Cref{sec:physics} exhibits the physical
realisation. \Cref{sec:scope} states scope and limitations and concludes.
